%!TeX root = main.tex
%!TeX root = main.tex
%!encoding = UTF-8 Unicode

\section{Active Current-Shunt Architecture and Implementation}

\begin{figure}[ht]
	\centering
	\tikzsetnextfilename{Diss_CompensatedCurrentShunt}
	\import{./figures/5_CurrentProbe/}{"Diss_CompensatedCurrentShunt.pdf_tex"}
	\caption{Simplified schematic overview of active current shunt.}\label{fig:ESBActiveShunt}
\end{figure}

\begin{figure}[ht]
	\centering
	\import{./figures/5_CurrentProbe/}{"Diss_Pic_AShunt.pdf_tex"}
	\caption{Photograph of active current shunt.}\label{fig:ESBActiveShuntPhoto}
\end{figure}

Therefore, a replica of the active shunt presented in~\cite{shillaber_gigahertz_2020,shillaber_ultrafast_2022} was designed. 
Its simplified equivalent circuit is shown in \fref{fig:ESBActiveShunt}, and a photograph of the prototype is shown in \fref{fig:ESBActiveShuntPhoto}.
This design offers multiple advantages:
First, a small shunt resistance can be used while still maintaining a sufficiently high sense output due to the amplification of the operational amplifier.
Furthermore, the low-pass filter of the active shunt compensates the high-pass characteristic of the shunt resistance, thereby achieving a remarkably high bandwidth.
Lastly, the compensation operational amplifier of the active shunt produces a defined output signal, which can either be fed into a filter stage for high-pass EM noise measurements or measured with an optically isolated voltage probe, thus providing galvanic isolation.



In contrast to~\cite{shillaber_gigahertz_2020,shillaber_ultrafast_2022}, extensive effort was required to sufficiently shield the sensitive active amplification from capacitive coupling from the switching-node potential, which would otherwise contaminate the measurement.
Best results were achieved with an ``EMC shield'' made of a nickel-silver alloy, which attenuates both capacitive and inductive coupling.
In addition to the shield, the transmission line between the passive shunt and the compensation amplifier was routed in the two inner PCB layers and enclosed by top and bottom ground planes connected by a via fence.



\subsection{Passive Shunt Element and Layout Realization}


\begin{figure}[ht]
	\centering
	\tikzsetnextfilename{Diss_PassiveShuntPCB_3D_wSchematic}
	\import{./figures/5_CurrentProbe/}{"Diss_PassiveShuntPCB_3D_wSchematic.pdf_tex"}
	\caption{Design of passive part of the current shunt.}\label{fig:ESBPassiveShunt}
\end{figure}






\subsection{Active Compensation for Bandwidth Extension and Impedance Matching}



\begin{figure}[ht]
	\centering
	\tikzsetnextfilename{Diss_Compensation_Shunt_ESB}
	\import{./figures/5_CurrentProbe/}{"Diss_Compensation_Shunt_ESB.pdf_tex"}
	\caption{Design of compensation part of the active current shunt.}\label{fig:ESBCompensation}
\end{figure}

\subsection{Amplification and Filter Path Design}

\subsubsection*{Low-Pass}

\begin{figure}[ht]
	\centering
	\tikzsetnextfilename{Diss_Lowpass_Shunt_ESB}
	\import{./figures/5_CurrentProbe/}{"Diss_Lowpass_Shunt_ESB.pdf_tex"}
	\caption{Design of low-pass part of the active current shunt.}\label{fig:ESBLowpass}
\end{figure}

\subsubsection*{High-Pass}

\begin{figure}[ht]
	\centering
	\tikzsetnextfilename{Diss_Highpass_Shunt_ESB}
	\import{./figures/5_CurrentProbe/}{"Diss_Highpass_Shunt_ESB.pdf_tex"}
	\caption{Design of high-pass part of the active current shunt.}\label{fig:ESBHighpass}
\end{figure}